\documentclass[11pt]{article}
\usepackage[T1]{fontenc}
\usepackage[margin=1in]{geometry}
\usepackage{amsmath,amssymb,amsthm,array,booktabs}
\usepackage[hidelinks]{hyperref}
\newtheorem{theorem}{Theorem}
\newtheorem{corollary}[theorem]{Corollary}
\newtheorem{lemma}[theorem]{Lemma}
\newcommand{\E}{\mathcal E}
\title{Counterexamples to Conjectures 3 and 4 of Huang, Stoica, Xia and Zhong\\
\large arXiv:2006.06716v1, Section 5.2, printed page 32}
\author{Haobo Ma\\\small ChronoAI Pte Ltd, Singapore
\and Wenlin Zhang\\\small National University of Singapore\\\small ChronoAI Pte Ltd, Singapore}
\date{November 3, 2026}
\begin{document}
\maketitle
\begin{abstract}
Conjectures 3 and 4 in Section 5.2 of Huang, Stoica, Xia and Zhong,
arXiv:2006.06716v1, printed page 32, are refuted by positive nonconstant solutions of their
tree equations of motion. For every integer $q\geq2$ and every real
$r>0$, $r\ne1$, we construct a solution on the $(q+1)$-regular tree.
Its lengths vary geometrically along a bi-infinite spine and are constant
within each attached branch. Every branch vertex has equal incident
lengths, although the solution is globally nonconstant. The cubic case
$q=r=2$ gives branch lengths $2^n$ and spine lengths $(5/3)2^n$.
More generally, any positive solution on a nontrivial tree of maximum degree at most
$q+1$ extends to $T_q$ by attaching constant-length branches without
changing its original vertex ratios. Applied to the source's classified
line solutions, this gives arbitrary two-choice scale patterns on every
$T_q$, $q\geq2$. The geometric-spine solutions have nonconstant tree
curvature.
\end{abstract}

\section{The statements and equations}
Huang, Stoica, Xia and Zhong~\cite{HSXZ}, Section 5.2, Conjecture 4
(printed p.~32), assert constancy of every positive solution on $T_2$.
Their Conjecture 3, on the same page, asserts global constancy if even one
vertex has equal incident lengths. Here $T_q$ denotes the infinite
$(q+1)$-regular tree. We give explicit counterexamples to both statements.

For positive edge lengths $\ell_{ij}=\ell_{ji}$, put
\begin{equation}\label{eq:definitions}
 c_i=\sum_{j\sim i}\ell_{ij}^{-1},\qquad
 d_i=\sum_{j\sim i}\ell_{ij}^{-2},\qquad R_i=c_i/d_i.
\end{equation}
Every denominator is positive. On a tree the unique simple path between
adjacent vertices is their edge; deleting backtracking from any other
finite walk decreases its length. Thus their shortest-path edge length
$P_{\langle ij\rangle}$ equals $\ell_{ij}$. Equation (3.1) of~\cite{HSXZ}
is precisely
\begin{equation}\label{eq:motion}
 \E_{ij}:=\frac{R_i^2+R_j^2}{\ell_{ij}}-R_i-R_j=0.
\end{equation}
No boundary condition is imposed here.

\begin{theorem}\label{thm:family}
Fix an integer $q\geq2$, a scale $t>0$ and a real $r>0$ with $r\ne1$.
There is a positive nonconstant solution of \eqref{eq:motion} on $T_q$
with the following lengths. Set
\begin{equation}\label{eq:parameters}
 s_n=t r^n\quad(n\in\mathbb Z),\qquad
 a=\frac{r^2+1}{r+1}.
\end{equation}
Choose a bi-infinite spine $(v_n)_{n\in\mathbb Z}$. The edge $v_nv_{n+1}$
has length $a s_n$. At each $v_n$ attach $q-1$ disjoint rooted
$q$-ary trees; every edge in these branches, including each attachment
edge, has length $s_n$. The ratio $R_i$ equals $s_n$ at $v_n$ and
at every vertex of a branch attached there.
\end{theorem}

\begin{proof}
Start with the doubly infinite path. A branch attached to $v_n$ has
vertices indexed by finite words over $\{1,\ldots,q\}$; its empty word
is joined to $v_n$, and a word is joined to each of its $q$ one-letter
extensions. Use a separate copy for each of the $q-1$ branches.
Each spine vertex has two spine neighbors and $q-1$ branch neighbors.
Each branch vertex has one parent and $q$ children. The graph is
therefore simple, locally finite and $(q+1)$-regular. It is connected:
following parents reaches the spine in finitely many steps. Each branch
and the spine are acyclic, and a branch has exactly one attachment edge,
so no cycle can pass from a branch to the spine and return. Hence the
graph is a realization of $T_q$.

All lengths are positive and symmetric. At $v_n$, the incident lengths
are $a s_n/r$, $a s_n$, and $q-1$ copies of $s_n$. Consequently
\begin{align}
 c_{v_n}&=\frac1{s_n}\left(\frac{r+1}{a}+q-1\right),\label{eq:c}\\
 d_{v_n}&=\frac1{s_n^2}\left(\frac{r^2+1}{a^2}+q-1\right).\label{eq:d}
\end{align}
The two parenthesized factors are equal, because
$a(r+1)=r^2+1$. They are strictly positive, so $R_{v_n}=s_n$.
At a branch vertex all $q+1$ incident lengths equal $s_n$; its sums
are $(q+1)/s_n$ and $(q+1)/s_n^2$, giving $R_i=s_n$ as well.

For a branch or attachment edge, both endpoint ratios and the edge
length equal $s_n$. Its residual is $2s_n^2/s_n-2s_n=0$.
For a spine edge $v_nv_{n+1}$, the endpoint ratios are $s_n$ and
$r s_n$, whence
\[
 \E_{v_nv_{n+1}}
 =s_n\left(\frac{1+r^2}{a}-1-r\right)=0.
\]
These cases exhaust the edges, and the residual is symmetric in its
endpoints. Finally the attachment edge at $v_0$ has length $t$,
whereas $v_0v_1$ has length $at$. These differ: $a-1=r(r-1)/(r+1)$
is nonzero. Thus the solution is not constant.
\end{proof}

\begin{corollary}\label{cor:conjectures}
Conjecture 4 of~\cite{HSXZ} is false. Conjecture 3 is false on $T_q$
for every $q\geq2$, even when equal incident lengths are required at
every vertex outside a bi-infinite path.
\end{corollary}
\begin{proof}
Take $q=2$, $r=2$, $t=1$ in Theorem~\ref{thm:family} to refute
Conjecture 4. For any $q\geq2$, every branch vertex has equal incident
lengths, so the same theorem contradicts the conclusion of Conjecture 3.
There is at least one branch at every spine vertex.
\end{proof}

For the cubic example the incident lengths at $v_n$ are
$(5/6)2^n$, $(5/3)2^n$, and $2^n$. Direct substitution gives
$c_{v_n}=14/(5\cdot2^n)$ and $d_{v_n}=14/(5\cdot2^{2n})$.
The attachment edge at $v_0$ has length $1$ and the next spine edge
has length $5/3$.

\section{Extension of the source's line solutions}
The cancellation at branch vertices gives a general extension principle.
It applies to all the line solutions in Section 5.1 of~\cite{HSXZ},
not just to a single geometric progression.

\begin{lemma}\label{lem:extension}
Let $q\geq2$ be an integer and let $H$ be a connected locally finite tree with at least
one edge and degrees at most $q+1$. Every positive solution of
\eqref{eq:motion} on $H$ extends to a positive solution on $T_q$,
preserving all original edge lengths and vertex ratios. At $i\in H$,
attach $m_i=q+1-\deg_H(i)$ disjoint rooted $q$-ary trees and give every
edge of each such branch, including its attachment, length $R_i$.
\end{lemma}
\begin{proof}
The original tree and the attached trees are connected and acyclic,
and every vertex now has degree $q+1$, so the enlarged graph is $T_q$.
At an original vertex the new sums satisfy
\[
 c'_i=c_i+\frac{m_i}{R_i},\qquad
 d'_i=d_i+\frac{m_i}{R_i^2},\qquad c'_i=R_i d'_i,
\]
because $c_i=R_i d_i$. Thus $R'_i=R_i$. Every new vertex in a branch
attached to $i$ has all incident lengths $R_i$ and hence ratio $R_i$.
The original edge residuals are unchanged. Every new edge, including
each attachment, has residual $2R_i^2/R_i-2R_i=0$.
\end{proof}

\begin{corollary}\label{cor:patterns}
Fix $q\geq2$ and any sequence $(b_n)_{n\in\mathbb Z}$ with
$b_n\in\{2,3\}$. Choose $L_0>0$ and define positive line-edge lengths
$L_n$ by $L_{n-1}/L_n=b_n$. There is a solution on $T_q$ with these
spine-edge lengths. At every vertex outside the spine the incident
lengths are equal, although the solution is nonconstant. If both choices
occur, the spine lengths are not one geometric progression.
\end{corollary}
\begin{proof}
Theorem 10 of~\cite{HSXZ}, Section 5.1, printed pp.~28--29, states that
these line lengths solve \eqref{eq:motion}: their successive reciprocal
ratios are $2$ or $(2+1)/(2-1)=3$. They are strictly decreasing, so the
line solution is nonconstant. At $v_n$ its ratio is
\[
 R_n=\frac{L_{n-1}^{-1}+L_n^{-1}}{L_{n-1}^{-2}+L_n^{-2}}>0.
\]
Apply Lemma~\ref{lem:extension} to the line, attaching $q-1$ branches
at each vertex. Both the original lengths and nonconstancy are preserved.
Two different successive length ratios exclude a single geometric
progression.
\end{proof}

Thus the source's prior classification on $T_1$ also supplies solutions
in every higher degree, including even $q$. For example one may choose
any periodic or nonperiodic binary sequence for $(b_n)$; no periodicity
assumption enters the proof.

\begin{corollary}\label{cor:global}
For a positive solution on a connected locally finite tree with at least
one edge, the edge lengths are globally constant if and only if the
vertex ratios $R_i$ are globally constant.
\end{corollary}
\begin{proof}
If $R_i=s$ everywhere, positivity gives $s>0$, and
\eqref{eq:motion} gives $2s^2/\ell_{ij}-2s=0$, hence $\ell_{ij}=s$.
Conversely, if every edge has length $s>0$, then
$c_i=\deg(i)/s$, $d_i=\deg(i)/s^2$, and $R_i=s$.
\end{proof}
This is a sufficient global replacement for the failed local equal-star
condition of Conjecture 3. The equal-star condition at one vertex does
not force its ratio to agree with all other vertex ratios.

\section{Curvature and adjacent statements}
\begin{corollary}\label{cor:curvature}
Every nonconstant geometric-spine solution of Theorem~\ref{thm:family}
has nonconstant tree curvature. Put
$B=q-1+(r+1)^2/(r^2+1)$. The curvatures of an edge wholly inside a
branch and of its attachment edge are respectively
\[
 K_{\rm branch}=\frac4{q+1}-2,\qquad
 K_{\rm attach}=\frac2B+\frac2{q+1}-2,
\]
and $K_{\rm attach}>K_{\rm branch}$. In the cubic $r=2$ example these
values are $-2/3$ and $-13/21$.
\end{corollary}
\begin{proof}
Equation (4.10) of~\cite{HSXZ}, printed p.~13, gives
\[
 K_{ij}=\frac2{\ell_{ij}^2}\left(\frac1{d_i}+\frac1{d_j}\right)
          -\frac{R_i+R_j}{\ell_{ij}}.
\]
At a spine vertex $d_i=B/s_n^2$ by \eqref{eq:d}; at a branch vertex
$d_i=(q+1)/s_n^2$. Substitute $\ell_{ij}=R_i=R_j=s_n$ for the two
edge types. Since $r\ne1$, $(r+1)^2<2(r^2+1)$, so $B<q+1$.
Both edge types exist for $q\geq2$, and their curvatures differ by
$2/B-2/(q+1)>0$.
\end{proof}

The construction specializes at $r=1$ to the constant solution.
For $r\ne1$, the spine lengths have infimum zero and are unbounded
above, since $r^n$ has these properties over $n\in\mathbb Z$.
In particular the spine does not approach one constant at its two
ends, so this family does not satisfy the asymptotically constant
boundary condition in~\cite{HSXZ}, Definition 13 and Theorem 2.

Positive nonconstant solutions for odd $q$ already appear in the
December 30, 2016 preprint of Gubser et al.~\cite{G}, Section 5,
especially equations (97)--(99), published in 2017,
and in~\cite{HSXZ}, Section 5.2, Theorem 12. Those constructions use
two equally sized groups of incident edges. In the present family
the two spine lengths and the branch length at a spine vertex are
pairwise distinct when $r\ne1$: in addition to $a\ne1$, the equality
$a/r=1$ would imply $r=1$. Thus these examples do not satisfy the
half-half hypothesis of~\cite{HSXZ}, Theorems 11--12.

The source's Remark 12 exhibits non-half-half solutions for odd $q$;
Theorem~\ref{thm:family} supplies further such solutions for odd $q\geq3$
and positive nonconstant solutions for every $q\geq2$.
Corollary~\ref{cor:patterns} also realizes the scale-pattern flexibility
discussed in Remark 13. These are solutions of \eqref{eq:motion}.
Corollary~\ref{cor:curvature} distinguishes them from the
constant-curvature metrics on finite trees studied in~\cite{BCH,Cheng}.

\begin{center}
\small
\begin{tabular}{@{}p{0.34\textwidth}p{0.58\textwidth}@{}}
\toprule
Source item in~\cite{HSXZ} & Coverage\\
\midrule
Section 5.2, Conjecture 4 (p.~32) & Refuted-and-corrected: positive nonconstant solutions exist on $T_2$ (Corollary~\ref{cor:conjectures}).\\
Section 5.2, Conjecture 3 (p.~32) & Refuted-and-corrected for every $q\geq2$; global constant ratios suffice (Corollaries~\ref{cor:conjectures}, \ref{cor:global}).\\
Section 5.1, Theorem 9 and Corollary 2 (pp.~27--28) & Prior proof cited: equal incident lengths force constancy on $T_1$; nonconstant line solutions are monotone.\\
Section 5.1, Theorem 10 (pp.~28--29) & Prior proof cited; all its positive line solutions lift to $T_q$ (Lemma~\ref{lem:extension}); two-choice patterns are proved in Corollary~\ref{cor:patterns}.\\
Section 5, Lemma 15 (p.~27) & Prior proof cited: common scaling preserves solutions; the parameter $t$ realizes it.\\
Section 5.2, Theorems 11--12 (pp.~29--30) & Prior proof cited for half-half solutions; their classification does not cover our geometric-spine family.\\
Section 5.2, Remark 11 (p.~30) & Prior curvature formulas for geometric half-half solutions; our different, nonconstant curvature is proved in Corollary~\ref{cor:curvature}.\\
Section 5.2, Remarks 12--13 (pp.~31--32) & Proved extensions: non-half-half solutions for odd $q\geq3$, solutions also in even $q\geq2$, and arbitrary two-choice spine patterns; general scale-pattern classification not covered.\\
Section 3, Theorems 1--3, Remarks 4--5 and Corollary 1 & Not covered: their boundary hypotheses differ; the geometric spine has lengths tending to zero and infinity at opposite ends.\\
Section 4, Theorems 4--8, Conjectures 1--2 and Remarks 6--10 & Not covered: action bounds, lattice and complete-graph extremizers concern different hypotheses and functionals.\\
\bottomrule
\end{tabular}
\end{center}

\section{Verification scope}
The declaration \texttt{result :} $\neg\,\texttt{claim}$ in the pinned
\href{https://github.com/the-omega-institute/trureturing/blob/50f2a97bb5e704842f2c32c6095f68a0dc62c779/D5/S3/Geometry/DiscreteGravity/TreeEquationsNonconstantRefutation.lean}{\texttt{TreeEquationsNonconstantRefutation.lean}}
proves the cubic counterexample.
Here \texttt{claim} is the implication from positivity, symmetry and
vanishing edge residuals on a locally finite cubic tree to global
constancy. The declaration proves its negation by the $q=r=2$, $t=1$
witness. Its supporting lemmas establish the tree and degree conditions,
the two vertex ratios, all edge residuals, and nonconstancy.
Theorem~\ref{thm:family}, Lemma~\ref{lem:extension} and
Corollaries~\ref{cor:conjectures}--\ref{cor:curvature} have written
proofs above; the line classification used in Corollary~\ref{cor:patterns}
is the source's Theorem 10. They are not additional Lean declarations.
The accompanying \texttt{verify.py} performs exact rational checks on
stated finite parameter sets, neighborhoods, periodic two-choice line
patterns and their extensions, and the curvature formulas. These checks
illustrate the formulas rather than prove the infinite statements.

\begin{thebibliography}{9}
\bibitem{HSXZ}
A. Huang, B. Stoica, X. Xia and X. Zhong,
\emph{Bounds on the Ricci curvature and solutions to the Einstein
equations for weighted graphs},
\href{https://arxiv.org/abs/2006.06716v1}{arXiv:2006.06716v1} (2020).

\bibitem{G}
S. S. Gubser, M. Heydeman, C. Jepsen, M. Marcolli, S. Parikh,
I. Saberi, B. Stoica and B. Trundy,
\emph{Edge length dynamics on graphs with applications to $p$-adic AdS/CFT},
J. High Energy Phys. \textbf{2017}, 157 (2017),
\href{https://doi.org/10.1007/JHEP06(2017)157}{doi:10.1007/JHEP06(2017)157};
\href{https://arxiv.org/abs/1612.09580v1}{arXiv:1612.09580v1}.

\bibitem{BCH}
S. Bai, H. Cheng and B. Hua,
\emph{Discrete Einstein metrics on trees},
\href{https://arxiv.org/abs/2604.22449v3}{arXiv:2604.22449v3} (2026), Section 1.

\bibitem{Cheng}
H. Cheng,
\emph{Positive-Curvature Discrete Einstein Metrics on Trees},
\href{https://arxiv.org/abs/2605.20862v2}{arXiv:2605.20862v2} (2026), Section 1.
\end{thebibliography}

\noindent\small\copyright\ 2026 Haobo Ma and Wenlin Zhang.
Licensed under \href{https://creativecommons.org/licenses/by/4.0/}{Creative Commons Attribution 4.0 International (CC BY 4.0)}.
\end{document}
